\section{Problemática global en el manejo de residuos y la generación de emisiones de GEI}
La gestión de residuos representa uno de los principales desafíos globales con implicaciones directas en el cambio climático y la sostenibilidad ambiental, ya que una amplia variedad de sectores productivos e industriales genera grandes volúmenes de residuos orgánicos e inorgánicos, siendo estos una fuente significativa de emisiones de gases de efecto invernadero (GEI). Por ejemplo, en el sector agrícola, los desechos orgánicos producidos durante las actividades productivas liberan metano ($CH_4$) debido a su descomposición anaeróbica \parencite{FAO2021}, mientras que otros sectores, como la industria manufacturera, el comercio y los servicios urbanos, también contribuyen considerablemente a la generación de residuos y emisiones contaminantes. Según el Quinto Informe de Evaluación del \textcite{IPCC2014}, actividades relacionadas con diversos sectores productivos son responsables de aproximadamente el $24\%$ de las emisiones globales de GEI, agravadas por prácticas como la quema de residuos y el uso intensivo de productos químicos. Estas actividades no solo incrementan las emisiones de GEI, sino que también causan problemas como la contaminación del suelo, el aire y el agua, afectando negativamente la biodiversidad y la salud humana. En respuesta a esta problemática, se requieren soluciones innovadoras y sostenibles que permitan una gestión eficiente de residuos en múltiples sectores, reduciendo así sus impactos ambientales globales \parencite{ahmedAgricultureClimateChange2020}.

